\documentclass[11pt]{article}
\usepackage[a4paper,margin=27mm]{geometry}
\usepackage{amsmath,amssymb,mathtools}
\usepackage{fontspec,unicode-math}
\setmainfont{Latin Modern Roman}
\setmathfont{Latin Modern Math}
\usepackage{xurl}
\usepackage[hidelinks]{hyperref}
\newcommand{\sref}[2]{\hyperlink{source:#1:#2}{[S#2]}}
\newcommand{\cataloguescope}[1]{\paragraph{Mathematical question.}#1}
\hypersetup{pdftitle={UnsolvedMath Problem 9400174}}
\begin{document}
% UnsolvedMath ID 9400174; source catalogue code AMR-093-0174
\setcounter{section}{9400173}
\section{\texorpdfstring{Wikipedia number-theory item 174: For any given integers $k\geq 1, b\geq 2, c\neq 0$, with gcd(k, c) = 1 and gcd(b, c) = 1, are there…}{Wikipedia number-theory item 174: For any given integers k>= 1, b>= 2, cneq 0, with gcd(k, c) = 1 and gcd(b, c) = 1, are there…}}\label{9400174}
{\small\sffamily UnsolvedMath ID 9400174 · Number Theory}
\subsection{Definitions and mathematical statement}

\paragraph{Shared notation.}$\mathbb N=\{1,2,\ldots\}$, and $\mathbb Z,\mathbb Q,\mathbb R,\mathbb C$ denote the integers, rationals, reals and complex numbers. Any other conventions are specified locally.


Fix integers $k\ge1$, $b\ge2$ and $c\ne0$ with $\gcd(k,c)=\gcd(b,c)=1$, where greatest common divisors are positive. Put $d=\gcd(k+c,b-1)$ and $A_n=(kb^n+c)/d$ for $n\ge1$. Since $b\equiv1\pmod d$, every $A_n$ is an integer. A prime value means $A_n>1$ is prime. As written, the universal assertion has a counterexample: $k=78557$, $b=2$, $c=1$ meets all the hypotheses and has $d=1$, while every $A_n=78557\cdot2^n+1$ is composite because $78557$ is a Sierpi\'nski number. \sref{9400174}{3}
\paragraph{Statement recorded in the supplied source.}
For every triple $(k,b,c)$ satisfying the displayed hypotheses, are there infinitely many $n\ge1$ for which $A_n$ is prime? The stated hypotheses alone do not imply this conclusion, as the displayed Sierpi\'nski counterexample shows. \sref{9400174}{2}
\subsection{Short English statement}
Does every indicated exponential sequence have infinitely many prime values? The stated conditions allow a known sequence consisting entirely of composite values.
\subsection{Sources}
\begin{description}
\item[\hypertarget{source:9400174:1}{S1}] ULAM.AI and the UnsolvedMath Contributors, UnsolvedMath: A Curated Collection of Open Mathematics Problems, record 9400174 (AMR-093-0174). CC BY 4.0. \url{https://huggingface.co/datasets/ulamai/UnsolvedMath}.
\item[\hypertarget{source:9400174:2}{S2}] Wikipedia contributors, \emph{List of unsolved problems in mathematics}, revision 1366636767 (2026) Number theory, Wikipedia number-theory item 174: For any given integers $k\geq 1, b\geq 2, c\neq 0$, with gcd(k, c) = 1 and gcd(b, c) = 1, are there…. \url{https://en.wikipedia.org/w/index.php?title=List_of_unsolved_problems_in_mathematics&oldid=1366636767#Number_theory}.
\item[\hypertarget{source:9400174:3}{S3}] OEIS Foundation and contributors, \emph{The On-Line Encyclopedia of Integer Sequences}, A076336: (Provable) Sierpiński numbers: odd numbers n such that for all k >= 1 the numbers n*2\textasciicircum{}k + 1 are composite., official data snapshot \texttt{693b5499c9d2} (2026) Name, definition and comments. \url{https://raw.githubusercontent.com/oeis/oeisdata/693b5499c9d248f570dd2d979e8bd1d7eecf689b/seq/A076/A076336.seq}.
\end{description}
\label{end:9400174}
\end{document}
